\section{Relazioni tra classi di complessità}
Ora che non siamo piú limitati a MdT precise, con input e output e con un solo nastro, e che possiamo usare la notazione asintotica, possiamo confrontare le classi di complessitá simulando le macchine di Turing tra loro. \\

\subsection{Teorema: terminazione delle MdT deterministiche che operano in spazio limitato}
Se una MdT $M$ opera in spazio $g(n)$, allora esiste una MdT $M'$ che opera in spazio $O(g(n))$ e termina per ogni input entro $O(c^{g(n)})$ passi.
\begin{proof}
$M' = (M_1, M_2)$ é la concatenazione di due macchine, e utilizza un nastro aggiuntivo. \\
$M_1$ scrive sul nastro aggiuntivo una stringa che codifica la configurazione di $M$, fatta in questo modo: \\
\[[\triangleright, q, \sqcap, \dots, \sqcap, 1, \dots, 1]\]
\begin{itemize}
	\item Il primo simbolo codifica lo stato corrente della macchina $M$ utilizzando $|Q|$ nuovi simboli.
	\item I successivi $g(n)$ simboli $\sqcap$ codificano il contenuto del nastro di $M$, utilizzando i $|\Sigma|$ simboli dell'alfabeto di $M$.
	\item Gli ultimi $g(n)$ simboli codificano la posizione della testina di $M$, rappresentata in unario tra le sue $g(n)$ posizioni possibili.
\end{itemize}
La macchina $M_1$ é simile a $M_g$, perché rappresenta due volte $g(n)$ in unario: una prima volta con $\sqcap^{g(n)}$ e una seconda con $1^{g(n)}$. \\
Lo spazio utilizzato é $O(g(n))$. \\

$M_2$ legge la configurazione di $M$ e utilizza gli altri nastri per simularla. \\
Il nastro con la configurazione viene trattato come un contatore, che viene incrementato dopo ogni simulazione. \\
Occorre solo fare attenzione al fatto che la base non é sempre la stessa:
\begin{itemize}
	\item La prima cella del contatore puó assumere $|Q|$ valori.
	\item Le successive $g(n)$ celle del contatore possono assumere $|\Sigma|$ valori.
	\item Le ultime $g(n)$ celle del contatore sono in unario.
\end{itemize}
Se la simulazione di $M$ termina, allora anche $M_2$ termina. \\
Se invece la simulazione non termina, il contatore prima o poi esaurirá le configurazioni possibili: a quel punto, faremo terminare $M_2$ anche se $M$ non termina. \\
Il massimo numero di configurazioni possibili é $|Q| \cdot |\Sigma|^{g(n)} \cdot {g(n)} = O(c^{g(n)})$.
\end{proof}
Questo teorema ci mostra come forzare la terminazione di una MdT che opera in spazio limitato.
\subsubsection{Corollario: $\SPACE(g(n)) \subseteq \TIME(c^{g(n)})$}
\subsubsection{Corollario: $\PSPACE \subseteq \EXP$}

\subsection{Teorema: simulazione di una MdT deterministica con una MdT non deterministica}
\label{theorem:deterministic_to_non_deterministic}
\begin{proof}
Le macchine di Turing deterministiche sono un caso particolare di quelle non deterministiche, in cui la relazione di transizione $\Delta$ associa a ogni coppia in $Q \times \Sigma$ una sola tripla di $Q \times \Sigma \times \{\leftarrow, \rightarrow, -\}$.
\end{proof}
\subsubsection{Corollario: $\TIME(f(n)) \subseteq \NTIME(f(n))$}
\subsubsection{Corollario: $\SPACE(g(n)) \subseteq \NSPACE(g(n))$}
\subsubsection{Corollario: $\P \subseteq \NP$}
\subsubsection{Corollario: $\EXP \subseteq \NEXP$}
\subsubsection{Corollario: $\PSPACE \subseteq \NPSPACE$}

\subsection{Teorema: simulazione di una MdT non deterministica che opera in tempo limitato con una MdT deterministica}
\label{theorem:non_deterministic_to_deterministic}
Data una MdT non deterministica $N$ che decide $L$ e opera in tempo $f(n)$. \\
Allora esiste una MdT deterministica $M$ che decide $L$ e opera in tempo $O(c^{f(n)})$ e spazio $O(f(n))$.
\begin{figure}[h!]
	\caption{\small Albero di computazione di una MdT non deterministica.}
	\centering
	\label{fig:albero_computazione}
	\includegraphics[width=0.8\textwidth]{albero_computazione.png}
\end{figure}
\begin{proof}
	Simuliamo $N$ con una MdT deterministica $M$. \\
	Possiamo vedere la computazione di $N$ come un albero di altezza $t = O(f(n))$ (figura \ref{fig:albero_computazione}). \\
	Ogni passo corrisponde a una \textit{scelta} tra al piú $d$ possibilitá. \\
	Ogni sequenza di $t$ passi corrisponde a una sequenza di $t$ scelte. \\
	Codifichiamo quindi un run con una stringa di $t$ numeri da $1$ a $d$. \\
	$M$ tiene su un nastro una stringa che codifica un run, partendo da quella con tutti $1$. \\
	Per quel run, $M$ simula le transizioni che avrebbe eseguito $N$. \\
	Poi incrementa la stringa che codifica il run e ripete la simulazione. \\
	In questo modo esplora tutti i possibili run di $N$. \\
	$M$ ha bisogno di un altro nastro per conservare sempre l'input $x$ e altri $k$ nastri per simulare quelli di $N$. \\
	Se la simulazione di un run termina in $"yes"$, $M$ accetta. \\
	Se tutte le simulazioni terminano in $"no"$, $M$ rifiuta. \\
	
	Il numero di possibili run é $\sum_{i=1}^{f(n)} d^i $. \\	
	Il tempo di una simulazione é $f(n)$, cioé quello che impiega $N$. \\
	Tempo: $f(n) \cdot \sum_{i=1}^{f(n)} d^i = O(d^{f(n)})$. \\
	
	Da \ref{teorema:spazio_limitato_superiormente_dal_tempo} possiamo anche dire che ogni simulazione usa spazio $O(f(n))$. \\
	E ogni simulazione puó sovrascrivere lo spazio usato dalla precedente. \\
	Lo spazio aggiuntivo richiesto da $M$ é quello per enumerare i run, che é logaritmico sul numero di run possibili:
	$O(\log(d^{f(n)})) = O(f(n) \cdot \log(d)) = O(f(n))$. \\
	Spazio: $O(f(n))$.
\end{proof}
\subsubsection{Corollario: $\NTIME(f(n)) \subseteq \bigcup_{c>1} \TIME(c^{f(n)})$}
\subsubsection{Corollario: $\NTIME(f(n)) \subseteq \SPACE(f(n))$}
\subsubsection{Corollario: $\NP \subseteq \EXP$}


\subsection{Teorema: simulazione di una MdT non deterministica che opera in spazio limitato con una MdT deterministica}
\label{theorem:non_deterministic_space_to_deterministic_time}
Data una MdT non deterministica $N$ con input e output, a $k$ nastri,
che decide $L$ e opera in spazio $g(n)$. \\
Allora esiste una MdT deterministica $M$ 
che decide $L$ e opera in tempo $O(c^{\log(n) + g(n)})$ e spazio $O(g^2(n))$.
\begin{proof}
La dimostrazione utilizza l'algoritmo per $\REACHABILITY$ mostrato nel teorema di Savitch (\ref{theorem:savitch}). \\
Possiamo vedere la computazione di $N$ come un grafo. \\
Ogni configurazione di $N$ é un vertice del grafo, 
e gli archi sono le transizioni possibili tra configurazioni. \\

Poiché $N$ é con input e output, 
possiamo semplificare la rappresentazione di una sua configurazione
conservando, del primo nastro, solo la posizione della testina,
e ignorando l'ultimo nastro. \\
Una configurazione di $N$ quindi é del tipo:
\[
C = (q, t, w_2, u_2, \dots, w_{k-1}, u_{k-1})
\]
Dove $t$ é la posizione della testina sul primo nastro. \\

Allora il grafo $G_N = (V, E)$ é definito da:
\[
V = \{C \mid C \text{ é una configurazione di } N\}
\]
\[
E = \{(C_1, C_2) \in V^2 \mid C_1 \xrightarrow{N} C_2\}
\]
Chiamiamo $C_s$ la configurazione iniziale e $C_{"yes"}$ la configurazione finale. \\

Per ipotesi, ogni nastro $w_i u_i$ é lungo al piú $g(n)$. \\
Quindi, il numero di configurazioni diverse (cioé $|V|$) é dato da:
\begin{itemize}
	\item $|Q|$ stati.
	\item $n+2$ possibili posizioni della testina sul primo nastro (contando $\triangleright$ e il primo $\sqcup$ dopo $x$).
	\item $|\Sigma|^{g(n)}$ simboli per gli altri $k-1$ nastri. \\ 
	Cioé $|\Sigma|^{(k-1) \cdot g(n)}$ possibili combinazioni. \\
\end{itemize}
\[
|V| = |Q| \cdot (n+2) \cdot |\Sigma|^{(k-1) \cdot g(n)} = O(n \cdot c^{g(n)}) = O(c^{\log(n) + g(n)})
\]

$M$ decide $L$ in tempo $O(|V|)$. \\
\[
x \in L \iff \exists \text{ un run di } N \text{ che accetta } x \iff \exists \text{ un cammino da } C_s \text{ a } C_{"yes"} \text{ in } G_N
\]
Ma questo é un problema di $\REACHABILITY$, che sappiamo essere risolvibile in tempo polinomiale rispetto al numero di vertici del grafo (da \ref{problema:reachability}) 
e spazio $O(\log^2(|V|)) = O(\log^2 (c^{g(n)})) = O(g^2(n))$ (da \ref{theorem:savitch}). \\

Restano da calcolare il tempo e lo spazio richiesti per costruire $G_N$. \\

Non possiamo costruire esplicitamente la matrice di adiacenza di $G_N$:
lo spazio richiesto é $|V|^2 = O(c^{2(\log(n) + g(n))})$, cioé esponenziale su $n$.
Ma se usa spazio esponenziale, anche il tempo per occuparlo é esponenziale. \\

Per risparmiare spazio, possiamo calcolare a ogni passo di $\REACHABILITY$ 
se esiste un arco tra due configurazioni $C_1$ e $C_2$,
semplicemente simulando un passo di $N$, in tempo polinomiale.
\end{proof}
\subsubsection{Corollario: $\NSPACE(g(n)) \subseteq \TIME(c^{\log(n) + g(n)})$}
\subsubsection{Corollario: $\NPSPACE(g(n) \subseteq \SPACE(g^2(n))$}
\subsubsection{Corollario: $\NPSPACE \subseteq \EXP$}
\begin{proof}
$
\NPSPACE = \bigcup_{i>0} \NSPACE(n^i) \subseteq \bigcup_{i>0} \TIME(c^{\log(n) + n^i})$ \\
$\bigcup_{i>0} \TIME(c^{\log(n)} \cdot c^{n^i}) = \bigcup_{i>0} \TIME(n c^{n^i}) = \bigcup_{i>0} \TIME(c^{n^i}) = \EXP
$
\end{proof}
\subsubsection{Corollario: $\PSPACE = \NPSPACE$}
\begin{proof}
$\PSPACE \subseteq \NPSPACE$ da \ref{theorem:deterministic_to_non_deterministic} \\
Applicando $\NSPACE(g(n)) \subseteq \SPACE(g^2(n))$ con $g(n) = n^i$: \\
$\NPSPACE = \bigcup_{i>0} \NSPACE(n^i) \subseteq \bigcup_{i>0} \SPACE(n^{2i}) = \bigcup_{i>0} \SPACE(n^i) = \PSPACE$
\end{proof}

\subsection{Proposizione: chiusura delle classi di complessitá deterministiche rispetto al complemento}
Se $C$ é una classe di complessitá deterministica, allora $\co C = C$
\begin{proof}
Sia $L \in C$ e sia $M$ una macchina di Turing che decide $L$ e opera in tempo $f(n)$ e spazio $g(n)$. \\
Allora possiamo costruire una macchina $M'$ che decide $\overline{L}$ e opera in tempo $f(n)$ e spazio $g(n)$, semplicemente invertendo gli stati $"yes"$ e $"no"$  di $M$. \\
\end{proof}
\subsubsection{Corollario: $\co \TIME(f(n)) = \TIME(f(n))$}
\subsubsection{Corollario: $\co \SPACE(f(n)) = \SPACE(f(n))$}
\subsubsection{Corollario: $\co \P = \P$}
\subsubsection{Corollario: $\co \EXP = \EXP$}
\subsubsection{Corollario: $\co \L = \L$}
\subsubsection{Corollario: $\co \PSPACE = \PSPACE = \NPSPACE = \co \NPSPACE$}
\subsubsection{Perché non é possibile fare lo stesso ragionamento per le classi non deterministiche}
Con una macchina non deterministica $N$ non basta invertire gli stati terminali. \\
Quello che vorremmo é una macchina $N'$ che accetta $\overline{L}$ perché:
\begin{itemize}
\item $x \in \overline{L}$ sse $\exists$ un run di $N'$ che accetta $x$.
\item $x \notin \overline{L}$ sse $\forall$ run di $N'$, il run non accetta $x$.
\end{itemize}
Invertendo gli stati $"yes$ e $"no"$ di $N$ non si otterrebbe una macchina che accetta il complemento del linguaggio. \\ 
Invertire $"yes"$ e $"no"$ porta una macchina che accetta $x$ se un solo run di $N$ lo rifiutava. 
\textbf{Un input accettato da una macchina non deterministica puó essere rifiutato da un run}, e quindi accettato anche dalla nuova macchina. \\

Vedremo che é possibile dimostrare la chiusura delle classi di complessitá di spazio non deterministiche rispetto al complemento
(e lo abbiamo giá fatto per $\NPSPACE$). \\
Per quanto riguarda le classi di complessitá di tempo non deterministiche, invece, il problema è aperto.


\subsection{Domande aperte}
\begin{itemize}
    \item $\P = \NP$?
    \item $\NP = \co\NP$?
    \item $\NP = \PSPACE$?
    \item $\PSPACE = \EXP$?
\end{itemize}
